% TOP_PROBLEM: {"id": 2653, "title": "Kourovka Notebook Problem 21.144", "category_id": 17, "category": {"id": 17, "name": "group_theory", "display_name": "Group Theory", "description": "Problems about groups, group actions, representations, and related algebraic structures.", "slug": "group-theory", "order_index": 17, "created_at": "2026-07-31T15:26:25.671Z"}, "status": "open", "problem_number": "KOU-21.144", "set_id": 10}
% FIRST_POSTED: 2026-10-01
% ENTRY_KIND: statement_only
% UnsolvedMath ID 2653; source catalogue code KOU-21.144
% !TeX program = lualatex
% Definition and sources only; this file is not an LLM solution attempt.
\documentclass[11pt,a4paper]{article}
\usepackage[margin=25mm]{geometry}
\usepackage{fontspec}
\setmainfont{lmroman10-regular.otf}[BoldFont=lmroman10-bold.otf,
 ItalicFont=lmroman10-italic.otf,BoldItalicFont=lmroman10-bolditalic.otf]
\usepackage{amsmath,amssymb,mathtools}
\usepackage{unicode-math}
\setmathfont{latinmodern-math.otf}
\AtBeginDocument{\let\setminus\smallsetminus}
\usepackage{xurl,hyperref}
\hypersetup{colorlinks=true,urlcolor=blue}
\setcounter{secnumdepth}{0}
\setlength{\parindent}{0pt}
\setlength{\parskip}{5pt}
\setlength{\emergencystretch}{3em}
\begin{document}
\section*{Kourovka Notebook Problem 21.144}\label{2653}
{\small UnsolvedMath ID 2653 · KOU-21.144 · Group Theory · Kourovka Notebook - New Problems, Issue 21}
\subsection{Definitions and mathematical statement}
Let $F$ be Thompson's group of orientation-preserving piecewise-linear homeomorphisms of $[0,1]$ under composition. Each map has finitely many breakpoints, each breakpoint is a dyadic rational $a/2^b$, and every slope is $2^k$ for some integer $k$.

The elementary amenable groups form the smallest class $\mathcal E$ containing all finite groups and all abelian groups and closed under taking subgroups, quotient groups, extensions, and directed unions. Closure under extensions means that $N\trianglelefteq G$, $N\in\mathcal E$, and $G/N\in\mathcal E$ imply $G\in\mathcal E$. A directed union is the union of a family of subgroups in which any two are contained in a third; membership of every group in that family implies membership of the union.

The conjecture is that every subgroup $H\le F$ satisfies
\[
 H\in\mathcal E\qquad\text{or}\qquad
 \exists\text{ an injective homomorphism }F\longrightarrow H.
\]
Equivalently, any subgroup of $F$ that is not elementary amenable must itself contain a copy of $F$.

Subgroups $H$ are arbitrary and need not be finitely generated or finitely presented. Containing a copy means abstract group containment through an injection, not merely having $F$ as a quotient or containing particular standard generators. The elementary amenable class is specified by its algebraic closure operations, so the first alternative makes a definite assertion stronger than merely belonging to some unspecified class of groups with invariant means.
\subsection{Short English statement}
Must every subgroup of Thompson’s group $F$ either be elementary amenable or contain a subgroup isomorphic to $F$?
\subsection{Sources}
\begin{itemize}
\item[S1] ULAM.AI and the UnsolvedMath Contributors, supplied problems.json, record 2653 (KOU-21.144), Kourovka Notebook - New Problems, Issue 21. Catalogue identity and source transcription; CC BY 4.0.
\url{https://huggingface.co/datasets/ulamai/UnsolvedMath}.
\item[S2] The Kourovka Notebook, 21st edition, version 47 (30 September 2026), new problems, Problem 21.144. Expanded formulation of the supplied catalogue statement.
\url{https://arxiv.org/pdf/1401.0300v47}.
\end{itemize}
\end{document}
